% !PRINTED 572
\weberpage{572}

\[
\psi_{\lambda,\mu}(\varepsilon)\,\psi_{\lambda,\mu}(\varepsilon^{-1})
= \frac{(\varepsilon^\lambda,\,r)(\varepsilon^{-\lambda},\,r)(\varepsilon^\mu,\,r)(\varepsilon^{-\mu},\,r)}
       {(\varepsilon^{\lambda+\mu},\,r)(\varepsilon^{-\lambda-\mu},\,r)}
\]
bildet. Und mit H\"ulfe von \S.\,169, (3) findet man so die Formel
\weq{5}{\psi_{\lambda,\mu}(\varepsilon)\,\psi_{\lambda,\mu}(\varepsilon^{-1}) = n,}
die nun wieder durch die Substitution $\lambda f, f$ f\"ur $\lambda, \mu$ specialisirt wird:
\weq{6}{\psi_\lambda(\alpha)\,\psi_\lambda(\alpha^{-1}) = n.}
Schliesslich wollen wir noch eine Eigenschaft der Functionen $\psi$ erw\"ahnen, die sich auch unmittelbar aus ihrer Definition \S.\,169, (5) ablesen l\"asst. Danach ist n\"amlich, wenn $\lambda, \mu, \nu$ drei solche Zahlen sind, dass $\lambda + \mu$, $\lambda + \mu + \nu$ nicht durch $m$ theilbar ist
\weq{7}{\psi_{\lambda,\mu}(\varepsilon)\,\psi_{\lambda+\mu,\nu}(\varepsilon) = \frac{(\varepsilon^\lambda,\,r)(\varepsilon^\mu,\,r)(\varepsilon^\nu,\,r)}{(\varepsilon^{\lambda+\mu+\nu},\,r)},}
d.\,h. das Product $\psi_{\lambda,\mu}(\varepsilon)\,\psi_{\lambda+\mu,\nu}(\varepsilon)$ bleibt unge\"andert, wenn die drei Zahlen $\lambda, \mu, \nu$ beliebig vertauscht werden (falls auch $\lambda + \nu$ und $\mu + \nu$ durch $m$ nicht theilbar sind). Danach ergiebt sich z.\,B.
\[
\psi_{\lambda,\mu}(\varepsilon)\,\psi_{\lambda+\mu,\nu}(\varepsilon) = \psi_{\mu,\nu}(\varepsilon)\,\psi_{\mu+\nu,\lambda}(\varepsilon),
\]
und wenn man darin $\lambda, \mu, \nu$ durch $2\lambda f, f, f$ ersetzt, und beachtet, dass $\psi_{2f,2\lambda f}(\varepsilon) = \psi_\lambda(\alpha^2)$ ist,
\weq{8}{\psi_{2\lambda}(\alpha)\,\psi_{2\lambda+1}(\alpha) = \psi_1(\alpha)\,\psi_\lambda(\alpha^2),}
vorausgesetzt, dass $2, 2\lambda + 1, 2\lambda + 2$ nicht durch $e$ theilbar sind.

\gesperrt{Jacobi} hat den Versuch gemacht, mit H\"ulfe dieser Formeln f\"ur ein gegebenes, als Primzahl vorausgesetztes $e$ die s\"ammtlichen $e - 2$ Functionen $\psi_\lambda(\alpha)$ durch die conjugirten Werthe einer einzigen, d.\,h. durch $\psi_1(\alpha), \psi_1(\alpha^2), \psi_1(\alpha^3) \ldots$ auszudr\"ucken, und hat die Rechnung bis $e = 23$ durchgef\"uhrt.

Seine Vermuthung aber, dass eine solche Darstellung allgemein m\"oglich sei, ist nach den Untersuchungen von \gesperrt{Kronecker} wahrscheinlich nicht richtig.

Wir wollen als Beispiel den Fall $e = 7$ durchf\"uhren. F\"ur diesen Fall ist
\[
\begin{aligned}
\lambda &= 1, 2, 3, 4, 5\\
\lambda' &= 1, 4, 5, 2, 3\\
\lambda'' &= 5, 4, 3, 2, 1,
\end{aligned}
\]
also ergeben die Formeln (2) und (4)
\[
\begin{aligned}
&\psi_4(\alpha) = \psi_2(\alpha^4), \quad \psi_5(\alpha) = \psi_3(\alpha^5)\\
&\psi_1(\alpha) = \psi_5(\alpha), \quad \psi_2(\alpha) = \psi_4(\alpha),
\end{aligned}
\]
und daraus noch
%CONT>>
